\scanpage{20}
% L01–04
\deleted{3) $\beta\ne0$, \inserted{inv.} $\deleted{\unsure{\beta|\rho}}p=2$ \inserted{et car $2$} (i.e. $U_\beta\simeq G_m$) Alors}

\[
\deleted{\mathrm{Aut}(X,\varphi)=\{1\}\qquad\uline{\mathrm{Aut}}(\unsure{X|G|Q},\varphi)\simeq{}_2U_\beta\overset{\text{can}}{\simeq}\mu_2}
\]
\deleted{On conclut à}

\uline{Proposition} $\uline{\mathrm{Aut}}(G,\varphi)\simeq {}_2(U_\beta)$

% L05–11
\uline{Corollaire 1} Supposons $\beta$ \inserted{$=2-\alpha$} inv, i.e. $U_\beta\overset{\text{loc}}{\simeq}G_m$, i.e.

\deleted{a) Si $\beta$ inv., et $2.1_S$ inversible, ou nul, alors}

\[
U_\beta\xrightarrow[\text{can}]{\sim}G_m,\qquad \inserted{\lambda\mapsto1+\lambda\beta}
\]
Alors : a) l'image $T$ de $U_\beta$ dans $G$ est l'unique tore maximal de $G$ \unsure{contnt|contenant} $u$. \deleted{Si \unsure{alors}}

b) $\uline{\mathrm{Aut}}(G,\varphi)={}_2T\overset{\text{can}}{\simeq}\mu_2$

\deleted{b) Si $2.1_S$ inv. ou $2.1_S=0$, Alors}

% L12–19
\uline{Corollaire 2} Supposons $\beta=2-\alpha=0$ i.e. $\alpha=2$, i.e. $U_\beta\overset{\text{loc}}{\simeq}G_a$ i.e. $\deleted{\unsure{U_\beta\ G.|U_\rho\simeq}}\ U_\beta\lhook\joinrel\longrightarrow\deleted{\unsure{G_m}}G_a$ \deleted{\unsure{$\lambda\mapsto\lambda$|$\lambda\mapsto1$}} est \unsure{multiplicatif|multiplicat.if}. Alors a) L'image $U$ de $U_\beta$ dans $G$ est l'unique \deleted{\unsure{ssgpe $\simeq G_a$}} ss groupe unipotent maximal de $G$ contnt $u$.

b) Si $2.1_S$ inv, Alors $\uline{\mathrm{Aut}}(G,\varphi)=\{1\}$

c) Si $2.1_S=0$ i.e. $S$ de car. $2$, Alors
\[
\uline{\mathrm{Aut}}(G,\varphi)=U\simeq G_a.
\]

% L20–24
\uline{Cor. 3} Soit $\gamma$ un gpe, et soit $\Gamma=\mathbb D_\infty\times\gamma$. Considérons la cat. $C(\gamma)$ des couples $(G,\varphi)$, $G$ forme de $\mathrm{GP}(1)_S$, $\varphi:\Gamma\longrightarrow G(S)$. Cette catégorie est équivalente à celle
% archive sheet number: 19
